\documentclass[10pt,a4paper]{amsart}
\usepackage[margin=19mm]{geometry}
\usepackage{amsmath,amssymb,mathtools}
\usepackage[scr=rsfs,cal=euler]{mathalfa}
\usepackage{xfrac}
\usepackage[unicode,hidelinks,bookmarks=false]{hyperref}
\newcommand{\ie}{i.e.\ }
\newcommand{\cf}{cf.\ }
\newcommand{\resp}{respectively\ }
\newcommand{\Subsection}{\subsection}
\providecommand{\nobreakdash}{\nobreak\hbox{-}}
\setlength{\emergencystretch}{2em}
\multlinegap0pt
\usepackage{amssymb}
\newtheorem{thm}{Theorem}
\newtheorem{pf}{Pf}
\title{Landau-Bloch type theorem for elliptic and $K$-quasiregular harmonic mappings}
\author{Vasudevarao Allu, Rohit Kumar}
\date{Source: arXiv:2212.09021v2 (CC BY 4.0)}
\begin{document}
\maketitle

\noindent\emph{Source and licence.} Landau-Bloch type theorem for elliptic and $K$-quasiregular harmonic mappings. Version arXiv:2212.09021v2 (CC BY 4.0), distributed by arXiv under the Creative Commons Attribution 4.0 International licence, http://creativecommons.org/licenses/by/4.0/, as recorded in the arXiv licence metadata for this exact version and shown on the abstract page https://arxiv.org/abs/2212.09021v2. Full licence notice: \texttt{licenses/arxiv-2212.09021-CC-BY-4.0.txt}.\par\smallskip

\noindent\textit{Editorial scope.} Counted unit: the article's thm Rohi-Vasu-P1-Theorem-001 (source0.tex lines 450--459). The article's complete original proof of it is delivered with it (source0.tex lines 521--603, one \texttt{proof} environment of 83 lines, which the article places away from the statement). No mathematics is written, repaired or re-derived: the statement and the proof are the article's own lines, in the article's own notation, and no retained line is substituted or re-flowed.  No further source-local context is needed: every cross-reference of every retained line resolves inside the two delivered spans.  Nothing was omitted from any retained span, so the package carries no omission record.  No retained line contains a citation, so the wrapper carries no bibliography and no bibliography entry.  No retained line contains a figure environment, an image-inclusion command or drawing code, so no image is delivered and the package declares no asset dependency.  Editorial formatting only, in addition: an AMS \texttt{amsart} preamble that restates the article's own declarations and macros used by the retained lines, namely the environments \texttt{thm}, \texttt{pf} on separate counters, together with the standard packages those lines need.  The retained environments print in this document as Theorem 1, Pf 1 in order of appearance, whereas the article prints the same items with the numbering of its own full document.  The complete native source of the article as distributed in the arXiv e-print for this exact version is bundled unchanged as \texttt{sources/135014/source0.tex}.
\begin{thm} \label{Rohi-Vasu-P1-Theorem-001}
 Let $f=h+\bar{g}$ be  $(K, K')$-elliptic harmonic mapping defined on unit disc $\mathbb{D}$, where 
 \begin{equation}\label{Rohi-Vasu-P1-eq-01}
 h(z)=\sum_{n=1}^{\infty} a_n z^n \quad \text{and} \quad g(z)=\sum_{n=1}^{\infty} b_n z^n.
 \end{equation}
If $\lambda_f(z) \leq \lambda$ for $z \in \mathbb{D}$, then
\begin{equation}\label{Rohi-Vasu-P1-eq-02}
\left|a_n\right|+\left|b_n\right| \leq \frac{K \lambda+\sqrt{K'}}{n}.
\end{equation}
\end{thm} 

\begin{pf} [{\bf Proof of Theorem   \ref{Rohi-Vasu-P1-Theorem-001}}]
Let $f=h+\bar{g}$ be $(K,K')$- elliptic harmonic mapping defined on the unit disc $\mathbb{D}$, where $h$ and $g$ are given by (\ref{Rohi-Vasu-P1-eq-01}). For fixed $r \in (0,1)$ and $z=r e^{i\theta}$ for some $\theta \in [0,2\pi]$, we have 
\begin{equation}\label{Rohi-Vasu-P1-eq-03}
\frac{\partial h}{\partial r}(r e^{i \theta})=\sum_{n=1}^{\infty} n a_n r^{n-1} e^{i n \theta} \ 
\text { and } \ \frac{\partial g}{\partial r}(r e^{i \theta})=\sum_{n=1}^{\infty} n b_n r^{n-1} e^{i n \theta}.
\end{equation}
Integrating both sides of (\ref{Rohi-Vasu-P1-eq-03}) with respect to $r$ from $ 0$ to $2\pi $ after multiplying $e^{-i n \theta}$ to the left equation and multiplying $e^{i n \theta}$ to the right equation, we obtain
$$
n a_n r^{n-1}=\frac{1}{2 \pi} \int_0^{2 \pi} \frac{\partial h( re^{i \theta}) }{\partial r} e^{-i n \theta} d \theta $$
and 
$$
n b_n r^{n-1}=\frac{1}{2 \pi} \int_0^{2 \pi} \frac{\partial g( re^{i \theta}) }{\partial r} e^{i n \theta} d \theta.  $$ 
for $n=1,2...$.
Therefore, we have
$$
 n\left|a_n\right| r^{n-1} \leq \frac{1}{2 \pi} \int_0^{2 \pi} \left| {\frac{\partial}{\partial r}}h(re^{i \theta}) \right|  d \theta  $$
 and 
 $$
  n\left|b_n\right| r^{n-1} \leq \frac{1}{2 \pi} \int_0^{2 \pi} \left| {\frac{\partial }{\partial r}}g(re^{i \theta}) \right|  d \theta.$$
  Hence, we get
  $$
 n\left(\left|a_n\right|+\left|b_n\right|\right) r^{n-1} \leq \frac{1}{2 \pi} \int_0^{2 \pi}\left(\left|\frac{\partial }{\partial r}h(re^{i \theta})\right|+\left|\frac{\partial}{\partial r} g(re^{i \theta})\right|\right) d \theta 
$$
It is easy to see that
$$
\begin{aligned}
\frac{\partial }{\partial r}h( re^{i \theta}) &=\frac{\partial h}{\partial z} \frac{\partial z}{\partial r}+\frac{\partial h}{\partial \bar{z}} \frac{\partial \bar{z}}{\partial r} \\
&=h_z e^{i \theta}+0 \quad \quad \ \  \left(\because h_{\bar{z}}=0\right)\\
&=h_z e^{i \theta}
\end{aligned}
$$
and similarly, $$
\frac{\partial}{\partial r} g( re^{i \theta})=h_z e^{i \theta}.
$$
We also see that
\begin{equation}\label{Rohi-Vasu-P1-eq-04}
\left|f_z\right|=\left|h_z\right| \text { and }\left|f_{\bar{z}}\right|=\left|g_z\right|.
\end{equation}
In view of (\ref{Rohi-Vasu-P1-eq-04}), we have

\begin{equation}\label{Rohi-Vasu-P1-Eq-001}
n\left(\left|a_n\right|+\left|b_n\right|\right) r^{n-1}  \leq \frac{1}{2 \pi} \int_0^{2 \pi}\left(\left|h_z e^{i \theta}\right|+\left|g_z e^{i \theta}\right|\right) d \theta 
=\frac{1}{2 \pi} \int_0^{2 \pi}\left(\left|f_z\right|+\left|f_{\bar{z}}\right|\right) d \theta. 
\end{equation}

 Since $\lambda_f(z) \leq \lambda$ and $f$ is sense preserving, we have
$$
\left|f_z\right| \leq \lambda+\left|f_{\bar{z}}\right|.$$
Using Lemma H, we obtain
$$
\left|f_z\right| \leq \lambda+k_1\left|f_z\right|+k_2,
$$
which implies that
$$
\left|f_z\right| \leq \frac{\lambda
+k_2}{1-k_1}.
$$
Therefore, we have 
$$
\begin{aligned}
|f_z|+|f_{\bar{z}}| & \leq |f_z|+k_1|f_z|+k_2 \\
&=\left(1+k_1\right)\left|f_z\right|+k_2 \\
& \leq\left(1+k_1\right) \frac{\lambda+k_2}{1-k_1}+k_2.
\end{aligned}
$$
In view of Lemma H, we have 
\begin{equation}\label{Rohi-Vasu-P1-Eq-002}
|f_z|+|f_{\bar{z}}| \leq K \lambda+\sqrt{K^{\prime}}.
\end{equation}

By equation (\ref{Rohi-Vasu-P1-Eq-001}) and equation (\ref{Rohi-Vasu-P1-Eq-002}), we obtain
$$
\begin{aligned}
n\left(\left|a_n\right|+\left|b_n\right|\right) r^{n-1} & \leq \frac{1}{2 \pi} \int_0^{2 \pi}\left(K \lambda+\sqrt{K^{\prime}}\right) d \theta \\
&=K \lambda+\sqrt{K^{\prime}}.
\end{aligned}
$$
Now by taking $r \rightarrow 1^{-}$, we obtain
$$
\left|a_n\right|+\left|b_n\right| \leq \frac{K \lambda+\sqrt{K^{\prime}}}{n}.
$$
This completes the theorem.
\end{pf}

\end{document}
